\documentclass{homework}
% \usepackage{graphicx} % Required for inserting images
\usepackage{amsmath}
\usepackage{gensymb}
\usepackage{subfig}
\usepackage{float}

\class{ECEN 5407}
\title{Homework 4}
\author{Sean Jennings}
\date{November 15, 2023}

\begin{document}
\maketitle

\newcommand{\dolperMWh}{\text{ \$/MWh}}
\newcommand{\perh}{\text{ /hr}}

\section{Problem 1}
\begin{letterenum}
\item Locational marginal pricing is the marginal price of electricity at a given bus a power
    system. Marginal price is the additional cost per unit of energy that results from
    the operating costs of generation (e.g. fuel prices). Constraints of the transmission
    system limit the amount of power that can be passed between buses, which means that
    prices between different buses may not fully equalize. This is called congestion.

\item The market clearing price is the price of electricity which meets the 
    demand of the system. Since the amount of power generated
    must be equal to the demand (which in the US, is generally inelastic to price),
    the market clearing price is determined by the marginal cost of power generation,
    which in realistic scenarios will vary as a function of generation/demand.
\item
    \begin{enumerate}
    
    \item Since the marginal cost at node 1 is less than that at node 2, it is
    cheapest to export as much power as possible to node 2. Since the capacity
    of the transmission line (500 MW) is less than the generation capacity at node 1
    (900 MW) minus the demand at node 1 (250 MW), the transmission line is the
    limiting factor. The flow on the line should therefore be 500 MW.

    \item Any additional power at node 1 is met entirely by the generation at node 1,
        so at node 1 the LMP is equal to the marginal cost of node 1 generation:
        \$20/MWh. Since the transmission line is operating at capacity, any additional
        power must be met by the generation at node 2. Therefore, the LMP at node
        2 is equal to the price of generation at node 2: \$45/MWh.

    \item We have 750 MW of generation at node 1 and 450 MW of generation at node 2
        so the total generation costs are:
        \begin{equation*}
            C_{gen} = 750 \MW \cdot 20\dolperMWh + 450 \MW \cdot 45\dolperMWh = \$35,250 \perh
        \end{equation*}

    \item The amount charged per unit of energy demanded at each node is the LMP,\
        which we will denote $\pi_1$ and $\pi_2$. The total charged at each node, $E_1$
        and $E_2$, is the demand times the LMP:
        \begin{align*}
            E_1 &= \pi_1 \cdot D_1 = 20\dolperMWh \cdot 250\MW = \$5,000 \perh \\
            E_2 &= \pi_2 \cdot D_2 = 45\dolperMWh \cdot 950\MW = \$42,750 \perh \\
        \end{align*}

    \item The profit of the system operator is the amount collected from payments
        of consumers minus the generation costs:
        \begin{equation*}
            E_1 + E_2 - C_{gen} = \$12,500  \perh
        \end{equation*}
    \end{enumerate}
\item 
    \begin{enumerate}
        \item Without the transmission constraint, the amount exported from node
            1 to node 2 is now limited by the generation capacity of generator 1.
            So generator 1 will operate at capacity of 900 MW. This results
            in 900 MW minus 250 MW = 650 MW exported to node 1 (and flowing over
            the line). Then, 300 MW
            is required from generator 2 to meet demand at node 2.

        \item Now, since generator 1 is running at capacity, Any additional
            load on the system must be met by generator 2. Transmission is no
            longer a constraint so the LMP at both nodes is equal to the
            marginal cost of generator 2: 45\$/MW.
    \end{enumerate}
\end{letterenum}

\section{Problem 2}

\begin{letterenum}
\item When the line is disconnected, the demand must be met entirely by local generation
    so that $P_A = D_A$ and $P_B = D_B$. Thus, the price of electricity at each
    node is:
    \begin{align*}
        MC_A = 20 + 0.2 \cdot 2000 = 420 \dolperMWh \\
        MC_B = 15 + 0.3 \cdot 1000 = 315 \dolperMWh \\
    \end{align*}
    The flow $F_{AB}$ is 0, of course.

\item Since electricity can be transmitted at will, $P_A$ and $P_B$ will be 
    determined such that the marginal cost of generation is equalized,
    and the total demand is met:
    \begin{align*}
        P_A + P_B &= D_T := D_A + D_B = 3000 \MW \\
        MC_A &= MC_B
    \end{align*}
    where $D_T$ is denotes the total demand. From these equations,
    we can first solve for the power output at A and B:
    \begin{align*}
        20 + 0.2 P_A &= 15 + 0.3 P_B \\
        5 + 0.2 (D_T - P_B) &= 0.3 P_B \\
        0.5 P_B &= 5 + 0.2 D_T \\
        P_B &= 2(5 + 0.2 \cdot 3000) \\
            &= 1210 \MW \\
        P_A = D_T - P_B = 1790 \MW.
    \end{align*}
    Then, the LMP at each node is same and equal to the marginal cost:
    \begin{equation*}
        LMP_A = LMP_B = MC_B = MC_A = (20 + 0.2 \cdot 1790) \dolperMWh = 378 \dolperMWh .
    \end{equation*}

    The line flow can be found via conservation of power at either node:
    \begin{align*}
        F_{AB} &= P_A - D_A = 1790 \MW - 2000 \MW = -210 \MW \\
               &= D_B - P_B = 1000 \MW - 1210 \MW = -210 \MW
    \end{align*}
    where $F_{AB} < 0$ indicates that power is flowing from node $B$ to node $A$.

\item Since the generation capacity at $A$ is less than the amount which equalizes
    the marginal cost, generator $A$ will run at capacity and the marginal price
    will determined by $MC_B$. With $P_A = 1000 \MW$, we have $P_B = 2000$ MW, so
    that:
    \begin{equation*}
        LMP_A = LMP_B = MC_B = 15 + 0.3 \cdot 2000 = 615 \dolperMWh.
    \end{equation*}
    The flow between the lines is:
    \begin{equation*}
        F_{AB} = P_A - D_A = -1000 \MW
    \end{equation*}

\item With the constraint on generator $B$, we now have generator $B$ running
    at capacity and the marginal cost determined by $A$:
    \begin{equation*}
     P_A = D_T - P_B = 3000 \MW - 900 \MW = 2100 \MW.
    \end{equation*}
    The LMPs are
    \begin{equation*}
        LMP_A = LMP_B = MC_A = 20 + 0.2 \cdot 2100 = 440 \dolperMWh
    \end{equation*}
    and the flow is
    \begin{equation*}
        F_{AB} = P_A - D_A = 2100 \MW - 2000 \MW = 100 \MW.
    \end{equation*}
    where $F_{AB} > 0$ now indicates that power is flowing from A to B.

    It makes sense that the LMP is lower than in part (c) since the marginal
    cost for generator $A$ increases at a slower rate than that of generator $B$
    (and we are far above the limit where the higher fixed cost is relevant).


\section{Problem 3}
% The problem does not specify whether the constraints given are constraints on
% active power or apparent power. Since reactances of the transmission lines
% are specified and "transmission losses" usually refers to losses due to resistivity
% rather than reactance, we will assume that reactive needs to be supplied. We will
% assume the constraints given are apparent power for generation and active
% power for transmission.

\begin{letterenum}
\item Let $P_1$ and $P_2$ be the power output of G1 and G2, respectively.
Since the marginal cost of G1 is less than that of G2, we first consider the
case where G2 is turned off entirely, i.e. $P_1 = 420 \MW$ and $P_2 = 0 \MW$.

Let $F_{ij}$ be the power flowing from bus $i$ to bus $j$, and $X_{ij}$ be the line impedance
between bus $i$ and bus $j$. Since no power is generated or consumed at bus 2, we have
$F_{12} = F_{23}$. We define $F_{123} := F_{12} = F_{23}$ as the power flowing
from bus 1 to 2 to 3, to make this equivalence
evident.

Per the figure, we also have
$X_{12} = X_{23} = X_{13} = 0.1$ p.u. Let $X_{123} = X_{12} + X_{23} = 0.2$ p.u.
be the sum of the series line impedances from bus 1 to 2 to 3.
The power is transmitted is divided between the 
transmission lines the power divider formula, which is derived from
Kirchoff's voltage and current laws:
\begin{align*}
    F_{123} &= \frac{X_{12}}{X_{123} + X_{13}} P_1 = \frac{0.1}{0.3} P_1 \\
        &= \frac{1}{3} P_1 \\
        &= \frac{1}{3} (420 \MW) = 140 \MW \\
    F_{13} &= \frac{X_{123}}{X_{123} + X_{13}} \\
        &= \frac{2}{3} P_1 \\
        &= \frac{2}{3} (420 \MW) = 280 \MW
\end{align*}

Since this violates the maximum power constraints for line 13, we must commit
to running G2. To minimize costs, we will run with G2 at minimum power:
$P_2 = 150$ MW. Then, $P_1$ is given by:
\begin{equation*}
    P_1 = D_T - P_2 = 420 \MW - 150 \MW = 270 \MW.
\end{equation*}
Now, the flow through transmission lines is:
\begin{align*}
    F_{123} &= \frac{1}{3} (270 \MW) = 90 \MW \\
    F_{13} &= \frac{2}{3} (270 \MW) = 180 \MW.
\end{align*}
Transmission lines constraints are no longer violated, so we have found the
optimal dispatch.

\item Since the system is running within transmission constraints and generation
constraints for G1, the electricity price (i.e.
marginal cost) is determined by the marginal cost of G1, the least expensive generator,
at \$20/MWh.

\item The total variable generation cost is given by:
\begin{align*}
    C_{gen} &= P_1 MC_1 + P_2 MC_2 \\
        &= 270 \MW \cdot 20 \dolperMWh + 150 \MW \cdot 35 \dolperMWh \\
        &= \$ 10,650 \perh
\end{align*}
\end{letterenum}

\section{Problem 4}
\begin{letterenum}
    \item The coined term that describes this diagram is "the duck curve."
        Between 2011 and 2016, the amount of distributed PV in the CAISO
        system increased significantly, resulting in a decrease
        in observed load during the day from the utility's perspective.
        However, a decrease in PV output (increase in load) when the sun goes down
        coincides with the evening peak in load (when people get home, start
        cooking dinner, and turn the lights on).

        This is an issue for grid operators because it increases the ramp
        rate of power generation. Since the ramp rate of any given generator is
        limited,
        this may require a larger number of generators to turn online 
        to obtain a sufficient ramping capacity. This translates
        to higher fixed costs for the grid as well as increased challenges
        in generator startup coordination for operators.

    \item Demand response is a good countermeasure because it can help reduce the
        overall evening peak, provide smoothing to the demand curve, and reduce
        the peak ramp rate. One DR service would be distributed battery
        energy storage. The grid could request that customers charge the battery during
        the middle of the day when PV output is high, and discharge during the evening
        peak, thus smoothing out the demand curve. This option would be particularly
        ideal because it does not require electricity consumers to change
        their behavior (e.g. not allowing people to turn on their microwave when
        they want to), apart from the initial capital investment.
        

        \indent Delayed EV charging and "smart thermostats" are two other DR services that
        would be particularly effective, given that each represent a substantial
        porportion of total demand (or projected future demand in the case of EV
        charging). These are somewhat less ideal than a battery dedicated
        to DR, since the modification of demand could conflict with
        consumer preferences: the EV might not be charged when the user wants to
        drive, or the house might be too warm/cold. Although more challenging to implement
        than the battery storage case,
        there is still a range DR curves which do not conflict with consumer preferences.
        For EV charging, many customers have generally predictable driving patterns
        (e.g. to and from work every day) and so charging during the day when
        they are at work and PV output is high, would likely not affect many people,
        most of the time. (However, when an unexpected trip does occur,
        the consumer would likely appreciate insight into their current state of
        EV charge and a convenient way to short-circuit the DR limitations.)
        
        \indent Similarly, there are a range of temperatures at which people are
        comfortable and thermal inertia can provide a degree of demand smoothing.
        So for instance, on a hot day,
        smart thermostats could front-load air-conditioning output during the day to drive the
        temperature down to the lower comfort limit, before the PV output drops.
        Then, the power requirement to maintain temperature under the upper
        temperature limit is lower later on, resulting in an overall
        decrease in peak evening demand. However, the tension between 
        smart thermostat DR and consumer
        comfort will impose limitations on the amount peak demand reduction
        that can be achieved.


\end{letterenum}
    






\end{letterenum}


\end{document}
